\chapter*{Conclusiones}
\addcontentsline{toc}{chapter}{Conclusiones}
La revisión de las fuentes consultadas muestra que el Renacimiento no fue un bloque homogéneo, sino un proceso articulado por la circulación de textos, imágenes y prácticas sociales. El humanismo proporcionó el marco conceptual para reordenar la herencia clásica, mientras que las artes visuales reforzaron nuevas formas de identidad individual y colectiva.\cite{Wind1958,WoodsMarsden1998}

Los avances en medicina y las respuestas a la sífilis evidencian la interacción entre conocimiento científico y gestión política de la población, confirmando que la salud pública fue un campo de experimentación administrativa.\cite{Arrizabalaga1997} La imprenta, por su parte, consolidó un espacio editorial europeo que normalizó géneros textuales, aceleró la crítica erudita y democratizó el acceso a los saberes.\cite{Pettegree2010}

Finalmente, las tensiones en torno a género y sexualidad revelan la fragilidad de las normas morales en sociedades en transformación. La combinación de estas dimensiones permite comprender el Renacimiento como una constelación de prácticas que redefinieron la relación entre tradición y modernidad en distintos espacios europeos.
